
\documentclass{article}
\usepackage{colortbl}
\usepackage{makecell}
\usepackage{multirow}
\usepackage{supertabular}

\begin{document}

\newcounter{utterance}

\centering \large Interaction Transcript for game `wordle\_withclue', experiment `medium\_frequency\_words\_clue\_no\_critic', episode 20 with Qwen3.5{-}4B.
\vspace{24pt}

{ \footnotesize  \setcounter{utterance}{1}
\setlength{\tabcolsep}{0pt}
\begin{supertabular}{c@{$\;$}|p{.15\linewidth}@{}p{.15\linewidth}p{.15\linewidth}p{.15\linewidth}p{.15\linewidth}p{.15\linewidth}}
    \# & \multicolumn{2}{c}{Player} && \multicolumn{2}{c}{Game Master} \\
    \hline

    \theutterance \stepcounter{utterance}  
    & & & \multicolumn{4}{p{0.6\linewidth}}{
        \cellcolor[rgb]{0.9,0.9,0.9}{
            \makecell[{{p{\linewidth}}}]{
                \texttt{\tiny{[P1$\langle$GM]}}
                \texttt{You are a language wizard who likes to guess words by using the given rules.} \\
\\ 
\texttt{Welcome to Wordle! You have six attempts to guess the target word, a valid English word of five lowercase letters (a{-}z). Please use the tags "explanation:" and "guess:" to provide a concise explanation for each guess.} \\
\\ 
\texttt{To help you make an informed guess, you will receive a clue for the word, such as} \\
\texttt{clue: snowy white.} \\
\\ 
\texttt{Here is an example guess based on the clue:} \\
\texttt{explanation: In the fairy tail Snow White, the girl is killed because she eats a poisoned apple. And the word apple has 5 letters.} \\
\texttt{guess: apple} \\
\\ 
\texttt{After each guess, your answer will be validated, and you will receive feedback indicating which letters are correct (green), which letters are correct but in the wrong position (yellow), and which letters are incorrect (red). This feedback can be useful in determining which letters to include or exclude in your next guess.} \\
\\ 
\texttt{For example, the feedback for "apple" might be:} \\
\texttt{guess\_feedback: a<yellow> p<yellow> p<green> l<yellow> e<red>} \\
\\ 
\texttt{The explanation should contain details about how both the guess\_feedback and clue work together in tandem to formulate a new guess.} \\
\\ 
\texttt{Do not generate any other text.} \\
\\ 
\texttt{Let's begin with your first guess.} \\
\\ 
\texttt{clue: winter wear} \\
            }
        }
    }
     \\ \\

    \theutterance \stepcounter{utterance}  
    & & \multicolumn{4}{p{0.6\linewidth}}{
        \cellcolor[rgb]{0.95,0.95,0.95}{
            \makecell[{{p{\linewidth}}}]{
                \texttt{\tiny{[P1$\leftrightarrow$P]}}
                \texttt{{[}DRAFT{]}} \\
\texttt{explanation: The clue "winter wear" suggests clothing suitable for cold weather. Common five{-}letter words fitting this description include "scarf" (too short), "boots" (plural, usually not the target), "coats" (plural), or "mitts" (plural). However, "mitts" is often pluralized as "mittens" (7 letters). Let's reconsider singular items. "Gloves" is 6 letters. "Socks" is 4. "Jackets" is 7. What about "coats"? No, plural. "Parka" is 5. "Parka" is a type of winter wear. Another option is "furs" (4). "Leather" is 7. "Down" is 4. "Fleece" is 6. "Therm" is 5 but not a word. "Woolen" is 6. "Knit" is 4. "Knits" is 5. "Knits" could be winter wear. "Sweats" is 6. "Sweat" is 5. "Sweat" is often associated with winter wear. "Parka" seems like a strong candidate. Let's try "parka".} \\
\texttt{guess: parka} \\
            }
        }
    }
    & \\ \\

    \theutterance \stepcounter{utterance}  
    & & \multicolumn{4}{p{0.6\linewidth}}{
        \cellcolor[rgb]{0.95,0.95,0.95}{
            \makecell[{{p{\linewidth}}}]{
                \texttt{\tiny{[P1$\leftrightarrow$P]}}
                \texttt{{[}SELF{-}FEEDBACK{]}} \\
\texttt{FORMAT: problem} \\
\texttt{RULES: problem} \\
\texttt{STRATEGY: problem} \\
\texttt{REASON: The response fails to use the required "explanation:" and "guess:" tags, violates the constraint of not generating extra text, and makes a weak guess based on an incorrect assumption that the target word must be singular.} \\
            }
        }
    }
    & \\ \\

    \theutterance \stepcounter{utterance}  
    & \multicolumn{4}{p{0.6\linewidth}}{
        \cellcolor[rgb]{0.561,0.737,0.561}{
            \makecell[{{p{\linewidth}}}]{
                \texttt{\tiny{[P1$\rangle$GM]}}
                \texttt{explanation: The clue "winter wear" suggests clothing for cold weather. While "parka" is a valid five{-}letter word, the previous guess assumed the target must be singular. Many winter wear items are commonly used in plural form in Wordle (e.g., "boots", "mittens" is 7, "gloves" is 6). "Socks" is 4 letters. "Scarves" is 6. "Coats" is 5 letters and fits the plural category of winter wear. "Parkas" is 6. "Snows" is 4. "Frosts" is 6. "Winters" is 7. "Heaters" is 7. "Warmers" is 7. "Thaw" is 4. "Chill" is 4. "Freeze" is 6. "Cold" is 4. "Snowy" is 5 but not wear. "Winter" is 6. "Coats" is a strong candidate for plural winter wear. Another option is "mitts" but that is 5 letters and often refers to mittens. Let's try "mitts" as it is 5 letters and fits the description of winter wear.} \\
\texttt{guess: mitts} \\
            }
        }
    }
    & & \\ \\

    \theutterance \stepcounter{utterance}  
    & & & \multicolumn{2}{p{0.3\linewidth}}{
        \cellcolor[rgb]{0.9,0.9,0.9}{
            \makecell[{{p{\linewidth}}}]{
                \texttt{\tiny{[GM$|$GM]}}
                \texttt{attempts = 1} \\
\texttt{target\_word = glove} \\
\texttt{target\_word\_clue = winter wear} \\
\texttt{guess = mitts} \\
\texttt{guess\_feedback = m<red> i<red> t<red> t<red> s<red>} \\
            }
        }
    }
    & & \\ \\

    \theutterance \stepcounter{utterance}  
    & & & \multicolumn{4}{p{0.6\linewidth}}{
        \cellcolor[rgb]{0.9,0.9,0.9}{
            \makecell[{{p{\linewidth}}}]{
                \texttt{\tiny{[P1$\langle$GM]}}
                \texttt{guess\_feedback: m<red> i<red> t<red> t<red> s<red>} \\
\\ 
\texttt{Provide your response only in this format.} \\
\texttt{explanation: details} \\
\texttt{guess: word} \\
            }
        }
    }
     \\ \\

    \theutterance \stepcounter{utterance}  
    & & \multicolumn{4}{p{0.6\linewidth}}{
        \cellcolor[rgb]{0.95,0.95,0.95}{
            \makecell[{{p{\linewidth}}}]{
                \texttt{\tiny{[P1$\leftrightarrow$P]}}
                \texttt{{[}DRAFT{]}} \\
\texttt{explanation: The previous guess "mitts" received all red feedback, indicating that the letters m, i, t, and s are not in the target word. This eliminates common winter wear items like mittens and socks. The clue "winter wear" still points to clothing, but without the letters m, i, t, or s. I need to find a five{-}letter word related to cold{-}weather clothing that uses different letters. "Parka" is a strong candidate as it is a quintessential winter coat, though it is singular. Other possibilities include "boots" (contains t and s, eliminated), "scarf" (5 letters, singular), or "coats" (contains t and s, eliminated). Since "scarf" is singular and fits the length, and "parka" is also singular, I will try "scarf" next.} \\
\texttt{guess: scarf} \\
            }
        }
    }
    & \\ \\

    \theutterance \stepcounter{utterance}  
    & & \multicolumn{4}{p{0.6\linewidth}}{
        \cellcolor[rgb]{0.95,0.95,0.95}{
            \makecell[{{p{\linewidth}}}]{
                \texttt{\tiny{[P1$\leftrightarrow$P]}}
                \texttt{{[}SELF{-}FEEDBACK{]}} \\
\texttt{FORMAT: ok} \\
\texttt{RULES: ok} \\
\texttt{STRATEGY: problem} \\
\texttt{REASON: The guess "scarf" is invalid because it contains the letter 'f', which was not present in the previous feedback, and the word is not a plural noun fitting the "winter wear" context better than other options.} \\
            }
        }
    }
    & \\ \\

    \theutterance \stepcounter{utterance}  
    & \multicolumn{4}{p{0.6\linewidth}}{
        \cellcolor[rgb]{0.561,0.737,0.561}{
            \makecell[{{p{\linewidth}}}]{
                \texttt{\tiny{[P1$\rangle$GM]}}
                \texttt{explanation: The previous guess "mitts" yielded all red feedback, eliminating m, i, t, and s. The clue "winter wear" suggests cold{-}weather clothing. "Boots" and "coats" are eliminated due to t and s. "Scarves" is six letters. "Parka" is singular. "Mitts" was plural but failed. "Scarves" is the most common plural winter wear item fitting the length and theme, despite the singular form "scarf" being an option. Since "mitts" failed, I will try "scarves" as a strong candidate for plural winter wear.} \\
\texttt{guess: scarves} \\
            }
        }
    }
    & & \\ \\

    \theutterance \stepcounter{utterance}  
    & & & \multicolumn{2}{p{0.3\linewidth}}{
        \cellcolor[rgb]{0.9,0.9,0.9}{
            \makecell[{{p{\linewidth}}}]{
                \texttt{\tiny{[GM$|$GM]}}
                \texttt{Error: The length of the guessed word is not 5.} \\
            }
        }
    }
    & & \\ \\

    \theutterance \stepcounter{utterance}  
    & & & \multicolumn{2}{p{0.3\linewidth}}{
        \cellcolor[rgb]{0.9,0.9,0.9}{
            \makecell[{{p{\linewidth}}}]{
                \texttt{\tiny{[GM$|$GM]}}
                \texttt{game\_result = ABORT} \\
            }
        }
    }
    & & \\ \\

\end{supertabular}
}

\end{document}
